% Paper 6 (GRPO-Registry) population section
\section{Populating the Registry from Real Artifacts}
\label{sec:p6-population}

A registry seeded with hypothetical entries would not test the artifact against
real reporting gaps. All \NSTACK{} seed stack records derive from concrete
artifacts of this program (Table~\ref{tab:p6-entries-summary}): framework config dumps
and runs for Qwen3-8B/GSM8K, the open Colab trainer audit, the managed-runtime
head-to-head, the N2 same-stack four-method run, two W\&B PPO/REINFORCE runs,
and the zvf-iter130 risk-index batch. Entries marked with a \texttt{run\_config}
block were regenerated from the run sources with per-field provenance; the
others remain hand-entered.

\begin{table}[t]
\centering
\caption{Seed stack entries by family (full per-entry table with badges and
spec coverage in Appendix~\ref{app:p6s-entries}). Gen.\ = entries regenerated
from run sources by \texttt{generate\_entries.py}; the rest are hand-entered.}
\label{tab:p6-entries-summary}
\small
\begin{tabular}{@{}lrlrl@{}}
\toprule
Family & $n$ & Labels & Gen. & Badge \\
\midrule
A: framework runs (Qwen3-8B, GSM8K) & 4 & grpo & 4 & 51--78 \\
B: open Colab trainer (E3, toy scale) & 4 & grpo, drgrpo, dapo & 4 & 84 \\
C: managed-runtime head-to-head (Qwen3.5-4B) & 4 & grpo, dapo, drgrpo, gspo & 4 & 78 \\
D: N2 same-stack variants (Qwen3.5-4B) & 3 & aero, gift, areal & 0 & 58--59 \\
W\&B PPO/REINFORCE runs & 2 & ppo-reinforce & 0 & 40 \\
E: zvf-iter130 risk-index fixtures & 11 & 9 methods + 2 tool-use & 0 & 36 \\
\midrule
Total & \NSTACK{} & & \NGEN{} & 36--84 \\
\bottomrule
\end{tabular}
\end{table}

\paragraph{Family A: the released config dumps.}
The four entries \texttt{trl/verl/openrlhf/tinker_grpo_qwen3-8b_gsm8k} are
mechanical translations of the 44-field per-framework dumps shipped at
\texttt{experiments/framework_config_dumps/}, since regenerated from the run
scripts where those exist. The dumps were written post hoc for reviewers and
do not describe the runs actually executed (e.g.\ placeholder checkpoint
hashes; TRL LoRA rank and completion length differ from the TRL script), so
the generated \texttt{run\_config} values take precedence. Translation preserved every honesty marker in the
dumps: the eleven \texttt{managed_by_tinker} nulls (including \texttt{kl.beta}
and \texttt{importance_sampling.granularity}), the genuine OpenRLHF deviations
($\beta{=}0.02$; sequence-level importance granularity), the Base-vs-Instruct
model confound on the managed runtime, and the fact that the veRL and OpenRLHF
rows of \texttt{framework_comparison.json} are dry-run placeholders (their
\texttt{telemetry} item is therefore null and their \texttt{outcomes} empty).

\paragraph{Family B: the open Colab trainer (E3 audit).}
Four entries record the arms of the E3 open-trainer audit (\ifdefined\ANONYMOUSBUILD
logged to W\&B; project name withheld for review\else W\&B project
\texttt{zvf-colab-experiments}\fi): base GRPO, Dr.~GRPO (both deltas
implemented), DAPO (clip-higher, dynamic sampling, token-level loss, and KL
removal implemented; overlong shaping recorded as \texttt{absent} because it
is inapplicable at toy completion lengths), and an adaptive-$G$ arm whose
\texttt{group_size_schedule} records the live escalation rule ($G$
$4{\to}6{\to}8$ on ZVF spikes). These runs are toy-scale and their outcome
fields (held-out delta $+0.50$ to $+0.575$; the DAPO arm's mean ZVF $0.00$ at
$+45\%$ rollouts) are directional evidence only---but the \emph{records} are
exact, because the trainer's backward pass is open.

\paragraph{Family C: the managed-runtime head-to-head.}
Four entries record the 12-cell head-to-head (Qwen3.5-4B, GSM8K, three seeds
per arm) in which the four variant labels spanned $0.034$ in last-10
training reward ($0.710$--$0.744$; not shown equivalent, and two seed-paired
contrasts exclude zero) while their delta status arrays differ
sharply from Family B's: every claimed delta on the closed stack is
\texttt{surrogate} or \texttt{unknown}, never fully \texttt{implemented}.

\paragraph{Family D: N2 same-stack isolated variants.}
Three entries (aero, gift, areal) record the same managed runtime as Family C
but with a different \emph{label} per cell. The N2 four-method run holds the
sampler, $G{=}8$, seed 0, prompt pool, and step count fixed---so the per-method
delta against GRPO is guaranteed to isolate the label, not the stack. Measured
last-10 reward/ZVF land in the same paired-bootstrap envelope as GRPO
\ifdefined\SUBMISSIONBUILD(Appendix~\ref{sec:p6-measured})\else(Section~\ref{sec:p6-measured})\fi; the entries make this measurable artifact
machine-queryable for the first time.

\paragraph{Family E: zvf-iter130 risk-index fixtures.}
Eleven records describe risk-index snapshots. They are not completed training
runs or validation of variant benefits. Their metadata is retained for
provenance, with current disclosure coverage in the per-entry table.

\paragraph{Snapshot warning.} Family badge ranges in this historical summary
are not a second scoring protocol; use the regenerated per-entry table for
current scores. The population includes dry-run placeholders and synthetic
risk-index fixtures, not 28 completed training runs. Family E cannot validate
variant performance or training ZVF.

\paragraph{What population taught us about our own reporting.}
The badge column is itself a finding. Our hand-entered config dumps---built for a
reviewer-requested fairness audit---scored low: they nail items 2, 3,
and 5 (KL, backend, group size) but expose item 1 only partially (a launcher
YAML records \texttt{importance_sampling.granularity}, not clip bounds or
advantage normalization) and items 6--7 not at all. The open Colab arms score
highest not because the runs are better but because an open backward pass plus an
explicit ``no KL, no decontamination'' declaration makes every field
\emph{reportable}. Reproducibility supplements of the kind commonly shipped
with RL papers, ours included, do not automatically satisfy MIN-REPORT.
